% GENERATED FILE. Edit canonical lessons, then run npm run generate:books.
% canonical-source-hash: fnv1a64:3e845330221e5515
% canonical-lessons: GE-C373-unvorhergesehen, GE-C373-spontan, GE-C373-beruecksichtigen, GE-C373-prognose, GE-C373-einplanen

\chapter{Mit allem rechnen}
\label{ch:ge-mit-allem-rechnen}
\hlchaptermodality{373}

\begin{chapteropening}
\textbf{By the end of this chapter:} \emph{I can talk about the unexpected and about spontaneous plans, forecasts, and what I allow for.}

Complete the last lesson of chapter 373: I can talk about the unexpected and about spontaneous plans, forecasts, and what I allow for.
\end{chapteropening}

\section[unvorhergesehen]{unvorhergesehen — unforeseen, unexpected}

\label{lesson:GE-C373-unvorhergesehen}



\begin{quote}
Before the new one: say the German for the schedule, then the German for to change your plans.
\end{quote}

\subsection*{You'll want to know: unvorhergesehen}
\textbf{unvorhergesehen} — \textquotedblleft{}unforeseen, unexpected\textquotedblright{}.

\begin{quote}
Wegen eines unvorhergesehenen Problems kommt der Techniker erst morgen.
\end{quote}

\begin{quote}
\emph{Because of an unforeseen problem the technician is not coming until tomorrow.}
\end{quote}

Built from \emph{un-}, \emph{vorher} — beforehand — and \emph{gesehen} — seen.

\begin{cousinweb}
\emph{unvorhergesehen} is built from \emph{un-} \textquotedblleft{}not\textquotedblright{}, \emph{vorher} \textquotedblleft{}beforehand\textquotedblright{} and \emph{gesehen}, the participle of \emph{sehen} \textquotedblleft{}to see\textquotedblright{}, Old High German \emph{sehan} (EtymDB): not seen in advance.
\end{cousinweb}

\begin{grammarlens}[title={What you've built}]
One more word for telling, explaining and describing at B1.
\end{grammarlens}

\subsection*{Your turn}
\begin{itemize}
\raggedright
  \item \emph{Say it:} \emph{unvorhergesehen}
  \item \emph{Say it:} \emph{unvorhergesehen}, once more
  \item \emph{Say it:} a sentence of your own with \emph{unvorhergesehen}
  \item \emph{Recall:} say the German for the schedule, then the German for to change your plans, then say \emph{unvorhergesehen} again
\end{itemize}

\begin{tcolorbox}[breakable,colback=teal!4,colframe=teal!35!black,title={Before you move on}]
What does unvorhergesehen mean? (\textquotedblleft{}Unforeseen, unexpected\textquotedblright{}.) Say it once more.
\end{tcolorbox}
\section[spontan]{spontan — spontaneous, on the spur of the moment}

\label{lesson:GE-C373-spontan}



\begin{quote}
Before the new one: say the German for to change your plans, then the German for unforeseen.
\end{quote}

\subsection*{You'll want to know: spontan}
\textbf{spontan} — \textquotedblleft{}spontaneous, on the spur of the moment\textquotedblright{}.

\begin{quote}
Wir sind ganz spontan ans Meer gefahren.
\end{quote}

\begin{quote}
\emph{We drove to the seaside on the spur of the moment.}
\end{quote}

Very common for plans made at the last minute: \emph{Ich bin da ganz spontan.}

\begin{cousinweb}
\emph{spontan} goes back to Latin \emph{spontāneus} \textquotedblleft{}of one's own accord\textquotedblright{} (EtymDB), from \emph{sponte} \textquotedblleft{}of one's free will\textquotedblright{} (Lewis \& Short).
\end{cousinweb}

\begin{grammarlens}[title={What you've built}]
One more word for telling, explaining and describing at B1.
\end{grammarlens}

\subsection*{Your turn}
\begin{itemize}
\raggedright
  \item \emph{Say it:} \emph{spontan}
  \item \emph{Say it:} \emph{spontan}, once more
  \item \emph{Say it:} a sentence of your own with \emph{spontan}
  \item \emph{Recall:} say the German for to change your plans, then the German for unforeseen, then say \emph{spontan} again
\end{itemize}

\begin{tcolorbox}[breakable,colback=teal!4,colframe=teal!35!black,title={Before you move on}]
What does spontan mean? (\textquotedblleft{}Spontaneous, on the spur of the moment\textquotedblright{}.) Say it once more.
\end{tcolorbox}
\section[berücksichtigen]{berücksichtigen — to take into account}

\label{lesson:GE-C373-beruecksichtigen}



\begin{quote}
Before the new one: say the German for unforeseen, then the German for spontaneous.
\end{quote}

\subsection*{You'll want to know: berücksichtigen}
\textbf{berücksichtigen} — \textquotedblleft{}to take into account\textquotedblright{}.

\begin{quote}
Bei der Reise müssen wir auch das Wetter berücksichtigen.
\end{quote}

\begin{quote}
\emph{For the trip we also have to take the weather into account.}
\end{quote}

Inseparable: \emph{er berücksichtigt, er hat berücksichtigt}.

\begin{cousinweb}
\emph{berücksichtigen} is built on \emph{Rücksicht} \textquotedblleft{}consideration\textquotedblright{}, literally a looking back, from \emph{Rücken} \textquotedblleft{}back\textquotedblright{} and \emph{Sicht} \textquotedblleft{}sight\textquotedblright{}, of the family of \emph{sehen} (EtymDB). \emph{Rücksicht} is usually explained as a rendering of Latin \emph{respectus} \textquotedblleft{}a looking back\textquotedblright{} (Lewis \& Short).
\end{cousinweb}

\begin{grammarlens}[title={What you've built}]
One more word for telling, explaining and describing at B1.
\end{grammarlens}

\subsection*{Your turn}
\begin{itemize}
\raggedright
  \item \emph{Say it:} \emph{berücksichtigen}
  \item \emph{Say it:} \emph{berücksichtigen}, once more
  \item \emph{Say it:} a sentence of your own with \emph{berücksichtigen}
  \item \emph{Recall:} say the German for unforeseen, then the German for spontaneous, then say \emph{berücksichtigen} again
\end{itemize}

\begin{tcolorbox}[breakable,colback=teal!4,colframe=teal!35!black,title={Before you move on}]
What does berücksichtigen mean? (\textquotedblleft{}To take into account\textquotedblright{}.) Say it once more.
\end{tcolorbox}
\section[die Prognose]{die Prognose — the forecast, the prognosis}

\label{lesson:GE-C373-prognose}



\begin{quote}
Before the new one: say the German for spontaneous, then the German for to take into account.
\end{quote}

\subsection*{You'll want to know: die Prognose}
\textbf{die Prognose} — \textquotedblleft{}the forecast, the prognosis\textquotedblright{}.

\begin{quote}
Die Prognose für das Wochenende ist gut: viel Sonne.
\end{quote}

\begin{quote}
\emph{The forecast for the weekend is good: lots of sun.}
\end{quote}

Plural: \emph{die Prognosen}. \emph{Eine Prognose stellen} — to make a forecast; a doctor's prognosis too.

\begin{cousinweb}
\emph{Prognose} is a learned borrowing from Late Latin \emph{prognōsis}, from Greek \emph{prógnōsis} \textquotedblleft{}foreknowledge\textquotedblright{}, built from \emph{pro} \textquotedblleft{}before\textquotedblright{} and \emph{gnôsis} \textquotedblleft{}knowing\textquotedblright{}. EtymDB has no entry for the German word.
\end{cousinweb}

\begin{grammarlens}[title={What you've built}]
One more word for telling, explaining and describing at B1.
\end{grammarlens}

\subsection*{Your turn}
\begin{itemize}
\raggedright
  \item \emph{Say it:} \emph{die Prognose}
  \item \emph{Say it:} \emph{die Prognose}, once more
  \item \emph{Say it:} a sentence of your own with \emph{die Prognose}
  \item \emph{Recall:} say the German for spontaneous, then the German for to take into account, then say \emph{die Prognose} again
\end{itemize}

\begin{tcolorbox}[breakable,colback=teal!4,colframe=teal!35!black,title={Before you move on}]
What does die Prognose mean? (\textquotedblleft{}The forecast, the prognosis\textquotedblright{}.) Say it once more.
\end{tcolorbox}
\section[einplanen]{einplanen — to allow for, to plan in}

\label{lesson:GE-C373-einplanen}



\begin{quote}
Before the new one: say the German for to take into account, then the German for the forecast.
\end{quote}

\subsection*{You'll want to know: einplanen}
\textbf{einplanen} — \textquotedblleft{}to allow for, to plan in\textquotedblright{}.

\begin{quote}
Für die Fahrt planen wir eine Stunde mehr ein.
\end{quote}

\begin{quote}
\emph{We are allowing an extra hour for the drive.}
\end{quote}

Separable: \emph{wir planen ein}, \emph{wir haben eingeplant}.

\begin{cousinweb}
\emph{einplanen} is built from \emph{ein} \textquotedblleft{}in\textquotedblright{} and \emph{planen} \textquotedblleft{}to plan\textquotedblright{}, from \emph{Plan}, borrowed from French \emph{plan} and going back to Latin \emph{plānus} \textquotedblleft{}flat\textquotedblright{} (EtymDB).
\end{cousinweb}

\begin{grammarlens}[title={What you've built}]
One more word for telling, explaining and describing at B1.
\end{grammarlens}

\subsection*{Your turn}
\begin{itemize}
\raggedright
  \item \emph{Say it:} \emph{einplanen}
  \item \emph{Say it:} \emph{einplanen}, once more
  \item \emph{Say it:} a sentence of your own with \emph{einplanen}
  \item \emph{Recall:} say the German for to take into account, then the German for the forecast, then say \emph{einplanen} again
\end{itemize}

\begin{tcolorbox}[breakable,colback=teal!4,colframe=teal!35!black,title={Before you move on}]
What does einplanen mean? (\textquotedblleft{}To allow for, to plan in\textquotedblright{}.) Say it once more.
\end{tcolorbox}
